\documentclass[10pt,a4paper]{amsart}
\usepackage[margin=19mm]{geometry}
\usepackage{amsmath,amssymb,mathtools}
\usepackage[scr=rsfs,cal=euler]{mathalfa}
\usepackage{xfrac}
\usepackage[unicode,hidelinks,bookmarks=false]{hyperref}
\newcommand{\ie}{i.e.\ }
\newcommand{\cf}{cf.\ }
\newcommand{\resp}{respectively\ }
\newcommand{\Subsection}{\subsection}
\providecommand{\nobreakdash}{\nobreak\hbox{-}}
\setlength{\emergencystretch}{2em}
\multlinegap0pt
\usepackage{bm}
\usepackage{bbm}
\usepackage{dsfont}
\usepackage{xspace}
\usepackage{cancel}
\usepackage{mathtools}
\usepackage{amsthm}
\makeatletter
\@ifundefined{thm}{\newtheorem{thm}{Theorem}}{}
\makeatother
\makeatletter
\expandafter\ifx\csname myeq\endcsname\relax
\newenvironment{myeq}[1][]
{\stepcounter{thm}\begin{equation}\tag{\thethm}{#1}}
{\end{equation}}
\fi
\expandafter\ifx\csname mysubsection\endcsname\relax
\newenvironment{mysubsection}[2][]
{\begin{subsec}\begin{upshape}\begin{bfseries}{#2.}
\end{bfseries}{#1}}
{\end{upshape}\end{subsec}}
\fi
\makeatother
\title[From $n$-Leibniz algebras and linear $n$-racks to the soluti]{From $n$-Leibniz algebras and linear $n$-racks to the solutions of the (higher analogue of) Yang-Baxter equation}
\author{Apurba Das, Suman Majhi}
\date{Source: arXiv:2508.07005v2 (CC BY 4.0)}
\begin{document}
\maketitle

\noindent\emph{Source and licence.} From $n$-Leibniz algebras and linear $n$-racks to the solutions of the (higher analogue of) Yang-Baxter equation. Version arXiv:2508.07005v2 (CC BY 4.0), distributed under the Creative Commons Attribution 4.0 International licence, as recorded in the arXiv licence metadata for this exact version and shown on the abstract page https://arxiv.org/abs/2508.07005v2. Full rights notice: licenses/arxiv-2508.07005-CC-BY-4.0.txt.\par\smallskip

\noindent\textit{Editorial scope.} Counted unit: the Thm whose source label is \texttt{thm-linn-lin} in the article (source0.tex lines 871--877, source label \texttt{thm-linn-lin}), delivered together with the article's own complete original proof of it (source0.tex lines 879--953, one \texttt{proof} environment of 75 lines). No mathematics is written, repaired or re-derived: statement and proof are the article's own text line for line. Required context retained uncounted: inv-prop (lines 787--789). Every definition, display and label of those lines is retained unchanged. Omitted from this package and not redistributed: 1--786, 790--870, 878--878, 954--2754. Nothing inside a retained excerpt is omitted or altered: every retained body is delivered exactly as the article's lines are written, so no omission receipt of any kind accompanies this package. Citations: the retained excerpts carry no citation command at all, so no bibliography entry of the article is transcribed and no reference is redistributed. Reference adaptation: none was needed and none was made; every reference inside the delivered unit resolves inside it, so no unresolved reference or citation is produced. Printed slips kept literal: no printed word, symbol, hypothesis or number of the counted statement or of its proof was altered. Layout and class: the article's own class is \texttt{amsart} with packages including epsf, graphicx, epsfig, amscd, xy, amsmath, latexsym, amssymb, amsthm, cite, textcomp, dsfont, setspace, lscape; the wrapper is the lane's pinned \texttt{amsart} editorial wrapper at 10pt on A4, which restates the theorem-like environments the retained text needs and adds the article's own macro definitions that the retained text needs (none), with extra wrapper packages (bm, bbm, dsfont, xspace, cancel, mathtools). The delivered numbering is the unit's own. Assets: no retained span contains a figure environment, a graphics-inclusion command, a PDF-page inclusion, a table or a TikZ picture; no image file is copied into this package, the package declares no asset dependency and no image file is redistributed with it. Licence: version arXiv:2508.07005v2 of this article is distributed under the Creative Commons Attribution 4.0 International licence, as recorded in the arXiv licence metadata for this exact version and shown on the abstract page https://arxiv.org/abs/2508.07005v2. A full rights notice is delivered at licenses/arxiv-2508.07005-CC-BY-4.0.txt. Characters: every retained excerpt is pure ASCII, so the delivered engine, XeLaTeX with its default Latin Modern text font, reports no missing character.
    \begin{align}\label{inv-prop}
        \ll \langle u, v_1^{(2)}, \ldots, v_{n-1}^{(2)} \rangle, v_{n-1}^{(1)}, \ldots, v_1^{(1)} \gg = \langle \ll u, v_1^{(2)}, \ldots, v_{n-1}^{(2)} \gg , v_{n-1}^{(1)}, \ldots, v_1^{(1)} \rangle = \varepsilon (v_1) \cdots \varepsilon (v_{n-1}) u.
    \end{align}

\begin{thm}\label{thm-linn-lin}
    Let $(C, \langle -, \ldots , - \rangle)$ be a linear $n$-rack where $C$ is cocommutative. Define a linear map $\triangleleft_{ \langle -, \ldots , - \rangle} : C^{\otimes (n-1)} \otimes C^{\otimes (n-1)} \rightarrow C^{\otimes (n-1)}$ by
    \begin{align}\label{first-n}
        (u_1 \otimes \cdots \otimes u_{n-1}) \triangleleft_{ \langle -, \ldots , - \rangle} (v_1 \otimes \cdots \otimes v_{n-1}) = \langle u_1 , v_1^{(1)}, \ldots, v_{n-1}^{(1)}   \rangle \otimes \cdots \otimes \langle u_{n-1} , v_1^{(n-1)}, \ldots, v_{n-1}^{(n-1)} \rangle,
    \end{align}
    for $u_1 \otimes \cdots \otimes u_{n-1} ~ \! ,  v_1 \otimes \cdots \otimes v_{n-1} \in C^{\otimes (n-1)}$. Then $(C^{\otimes (n-1)}, \triangleleft_{ \langle -, \ldots , - \rangle} ) $ is a linear rack. Moreover, if $f : (C, \langle -, \ldots , - \rangle) \rightarrow (C', \langle -, \ldots , - \rangle')$ is a homomorphism of cocommutative linear $n$-racks then the map $f^{\otimes (n-1)} : C^{\otimes (n-1)} \rightarrow C'^{\otimes (n-1)}$ is a homomorphism of linear racks from $(C^{\otimes (n-1)}, \triangleleft_{\langle -, \ldots , - \rangle })$ to $(C'^{\otimes (n-1)}, \triangleleft_{\langle -, \ldots , - \rangle' })$.
\end{thm}

\begin{proof} We divide the proof into the following steps.

  \noindent  {\sf Step 1.} {\em $\triangleleft_{\langle -, \ldots, - \rangle } $ is a coalgebra homomorphism.} For any $u_1 \otimes \cdots \otimes u_{n-1} ~\! , ~ \! v_1 \otimes \cdots \otimes v_{n-1} \in C^{\otimes (n-1)}$,
     \begin{align*}
       & \big( (\triangleleft_{\langle -, \ldots , - \rangle } \otimes \triangleleft_{\langle -, \ldots , - \rangle }) \circ \Delta_{C^{\otimes (n-1)} \otimes C^{\otimes (n-1)}} \big) ((u_1 \otimes \cdots \otimes u_{n-1}) \otimes (v_1 \otimes \cdots \otimes v_{n-1})) \\
        &= \big( (u_1^{(1)} \otimes \cdots \otimes u_{n-1}^{(1)}) \triangleleft_{\langle -, \ldots , - \rangle } (v_1^{(1)} \otimes \cdots \otimes v_{n-1}^{(1)}) \big) \otimes \big( (u_1^{(2)} \otimes \cdots \otimes u_{n-1}^{(2)}) \triangleleft_{\langle -, \ldots , - \rangle } (v_1^{(2)} \otimes \cdots \otimes v_{n-1}^{(2)}) \big) \\
        &= \big(  \langle u_1^{(1)} , v_1^{(1)(1)} , \ldots, v_{n-1}^{(1)(1)} \rangle \otimes \cdots \otimes \langle u_{n-1}^{(1)} , v_1^{(1)(n-1)} , \ldots, v_{n-1}^{(1)(n-1)} \rangle  \big) \\
        & \qquad \qquad \qquad \qquad \otimes \big(  \langle u_1^{(2)} , v_1^{(2)(1)} , \ldots, v_{n-1}^{(2)(1)} \rangle \otimes \cdots \otimes  \langle u_{n-1}^{(2)} , v_1^{(2)(n-1)} , \ldots, v_{n-1}^{(2)(n-1)} \rangle  \big) \\
        &= \big(   \langle u_1, v_1^{(1)}, \ldots, v_{n-1}^{(1)} \rangle^{(1)} \otimes \cdots \otimes \langle  u_{n-1}, v_1^{(n-1)}, \ldots, v_{n-1}^{(n-1)}  \rangle^{(1)} \big) \\
        & \qquad \qquad \qquad \qquad \otimes \big(   \langle u_1, v_1^{(1)}, \ldots, v_{n-1}^{(1)} \rangle^{(2)} \otimes \cdots \otimes \langle  u_{n-1}, v_1^{(n-1)}, \ldots, v_{n-1}^{(n-1)}  \rangle^{(2)}    \big) \\
        &= \Delta_{C^\otimes (n-1)} \big(   \langle u_1, v_1^{(1)}, \ldots, v_{n-1}^{(1)} \rangle \otimes \cdots \otimes  \langle u_{n-1}, v_1^{(n-1)}, \ldots, v_{n-1}^{(n-1)} \rangle \big) \\
        &= (\Delta_{C^\otimes (n-1)} \circ \triangleleft_{\langle -, \ldots , - \rangle }) ((u_1 \otimes \cdots \otimes u_{n-1}) \otimes (v_1 \otimes \cdots \otimes v_{n-1}))
     \end{align*}
     and
     \begin{align*}
         &(\varepsilon_{C^{\otimes (n-1)}} \circ \triangleleft_{\langle -, \ldots , - \rangle }) \big(   (u_1 \otimes \cdots \otimes u_{n-1}) \otimes (v_1 \otimes \cdots \otimes v_{n-1})   \big) \\
         &= \varepsilon_{C^{\otimes (n-1)}} \big(   \langle u_1, v_1^{(1)}, \ldots, v_{n-1}^{(1)} \rangle \otimes  \langle u_2, v_1^{(2)}, \ldots, v_{n-1}^{(2)} \rangle \otimes \cdots \otimes  \langle u_{n-1}, v_1^{(n-1)}, \ldots, v_{n-1}^{(n-1)} \rangle   \big) \\
         &= \varepsilon (u_1) \varepsilon (v_1^{(1)}) \cdots \varepsilon (v_{n-1}^{(1)}) 
         \varepsilon (u_2) \varepsilon (v_1^{(2)}) \cdots \varepsilon (v_{n-1}^{(2)}) \cdots \varepsilon (u_{n-1}) \varepsilon (v_1^{(n-1)}) \cdots \varepsilon (v_{n-1}^{(n-1)}) \\
         &= \varepsilon (u_1) \varepsilon (u_2) \cdots \varepsilon (u_{n-1}) \big(  \varepsilon (v_1^{(1)}) \varepsilon (v_1^{(2)}) \cdots \varepsilon (v_1^{(n-1)})    \big) \cdots \big(  \varepsilon (v_{n-1}^{(1)}) \varepsilon (v_{n-1}^{(2)}) \cdots \varepsilon (v_{n-1}^{(n-1)})   \big) \\
         &= \varepsilon (u_1) \varepsilon (u_2) \cdots \varepsilon (u_{n-1}) \varepsilon (v_1) \cdots \varepsilon (v_{n-1}) \\
         &= (\varepsilon_{C^{\otimes (n-1)} \otimes C^{\otimes (n-1)}}) \big(   (u_1 \otimes \cdots \otimes u_{n-1}) \otimes (v_1 \otimes \cdots \otimes v_{n-1})   \big).
     \end{align*}

     \medskip

   \noindent   {\sf Step 2.} { $\triangleleft_{\langle -, \ldots, - \rangle }$ {\em is a linear shelf.}} For any $u_1  \otimes \cdots \otimes u_{n-1}, ~ \! v_1  \otimes \cdots \otimes v_{n-1} ~ \! , ~ \! w_1 \otimes \cdots \otimes w_{n-1} \in C^{\otimes (n-1)}$, we have
   \begin{align*}
        &\big(  (u_1  \otimes \cdots \otimes u_{n-1}) \triangleleft_{\langle -, \ldots , - \rangle } ( v_1  \otimes \cdots \otimes v_{n-1}) \big) \triangleleft_{\langle -, \ldots , - \rangle } (w_1 \otimes \cdots \otimes w_{n-1}) \\
        &= \big( \langle u_1, v_1^{(1)}, \ldots, v_{n-1}^{(1)} \rangle \otimes \cdots \otimes \langle  u_{n-1}, v_1^{(n-1)}, \ldots, v_{n-1}^{(n-1)} \rangle   \big) \triangleleft_{\langle -, \ldots , - \rangle } (w_1 \otimes \cdots \otimes w_{n-1}) \\
        &= \big\langle \langle u_1, v_1^{(1)}, \ldots, v_{n-1}^{(1)} \rangle , w_1^{(1)}, \ldots, w_{n-1}^{(1)} \big\rangle \otimes \cdots \otimes  \big\langle \langle u_{n-1}, v_1^{(n-1)}, \ldots, v_{n-1}^{(n-1)} \rangle , w_1^{(n-1)}, \ldots, w_{n-1}^{(n-1)} \big\rangle \\
        &= \big\langle \langle u_1, w_1^{(1)(1)}, \ldots, w_{n-1}^{(1) (1)} \rangle, \langle 
  v_1^{(1)}, w_1^{(1) (2)}, \ldots, w_{n-1}^{(1) (2)}  \rangle, \ldots, \langle v_{n-1}^{(1)}, w_1^{(1) (n)} , \ldots, w_{n-1}^{(1)(n)} \rangle  \big\rangle  \otimes \cdots \otimes \\
  & \quad \big\langle \langle u_{n-1}, w_1^{(n-1)(1)}, \ldots, w_{n-1}^{(n-1) (1)} \rangle, \langle 
  v_1^{(n-1)}, w_1^{(n-1) (2)}, \ldots, w_{n-1}^{(n-1) (2)}  \rangle, \ldots, \langle v_{n-1}^{(n-1)}, w_1^{(n-1) (n)} , \ldots, w_{n-1}^{(n-1)(n)} \rangle  \big\rangle \\
  &= \big\langle \langle u_1, w_1^{(1)(1)}, \ldots, w_{n-1}^{(1)(1)} \rangle, \langle v_1, w_1^{(2)(1)}, \ldots, w_{n-1}^{(2)(1)}   \rangle^{(1)}, \ldots, \langle v_{n-1}, w_1^{(2)(n-1)}, \ldots, w_{n-1}^{(2)(n-1)}   \rangle^{(1)} \big\rangle  \otimes \cdots \otimes \\
  & \quad \big\langle \langle u_{n-1}, w_1^{(1)(n-1)}, \ldots, w_{n-1}^{(1)(n-1)} \rangle, \langle v_1, w_1^{(2)(1)}, \ldots, w_{n-1}^{(2)(1)}   \rangle^{(n-1)}, \ldots, \langle v_{n-1}, w_1^{(2)(n-1)}, \ldots, w_{n-1}^{(2)(n-1)}   \rangle^{(n-1)} \big\rangle \\
  &= \big( \langle u_1, w_1^{(1) (1)}, \ldots, w_{n-1}^{(1) (1)}    \rangle \otimes \cdots \otimes  \langle u_{n-1}, w_1^{ (1)(n-1)} , \ldots, w_{n-1}^{ (1)(n-1)} \otimes    \rangle    \big) \\
  & \qquad \qquad \quad \qquad \qquad \qquad \triangleleft_{\langle -, \ldots , - \rangle }   \big(  \langle v_1, w_1^{(2) (1)}, \ldots, w_{n-1}^{(2) (1)} \rangle \otimes \cdots \otimes  \langle v_{n-1}, w_1^{ (2)(n-1)} , \ldots, w_{n-1}^{ (2)(n-1)} \otimes  \rangle  \big) \\
  &= \big(  {\scriptstyle  (u_1 \otimes \cdots \otimes u_{n-1})} \triangleleft_{\langle -, \ldots , - \rangle } {\scriptstyle  (w_1^{(1)} \otimes \cdots \otimes w_{n-1}^{(1)})}   \big) \triangleleft_{\langle -, \ldots , - \rangle } \big( {\scriptstyle  (v_1 \otimes \cdots \otimes v_{n-1})} \triangleleft_{\langle -, \ldots , - \rangle } {\scriptstyle (w_1^{(2)} \otimes \cdots \otimes w_{n-1}^{(2)})}    \big).
   \end{align*}

   \medskip

  \noindent   {\sf Step 3.} {$ \triangleleft_{\langle -, \ldots, - \rangle } $ {\em is a linear rack.} } Since $(C, \langle -, \ldots, - \rangle)$ is a linear $n$-rack, there exists a coalgebra homomorphism $\ll - , \ldots , - \gg : C^{\otimes n} \rightarrow C$ that makes $(C, \ll - , \ldots , - \gg)$ into a linear $n$-shelf satisfying additionally (\ref{inv-prop}). We define a linear map $\triangleleft_{\ll - , \ldots , - \gg} : C^{\otimes (n-1)} \otimes C^{\otimes (n-1)} \rightarrow C^{\otimes (n-1)}$ by
  \begin{align*}
   {\scriptstyle (u_1 \otimes \cdots \otimes u_{n-1})} \triangleleft_{\ll - , \ldots , - \gg }  {\scriptstyle (v_1 \otimes \cdots \otimes v_{n-1}) }
   :=  \ll u_1, v_{n-1}^{(1)}, \ldots, v_1^{(1)} \gg \otimes \cdots \otimes \ll u_{n-1}, v_{n-1}^{(n-1)}, \ldots, v_1^{(n-1)} \gg.
  \end{align*}
  Then   $(C^{\otimes n-1}, \triangleleft_{ \ll - , \ldots, \gg})$ is a linear shelf. Additionally, 
     \begin{align*}
         &\big( (u_1 \otimes \cdots u_{n-1}) \triangleleft_{\langle -, \ldots, - \rangle} (v_1^{(2)} \otimes \cdots \otimes v_{n-1}^{(2)}) \big) \triangleleft_{ \ll - , \cdots ,- \gg } (v_1^{(1)} \otimes \cdots \otimes v_{n-1}^{(1)}) \\
         &= \big(   \langle u_1, v_1^{(2)(1)}, \ldots, v_{n-1}^{(2) (1)} \rangle \otimes \cdots \otimes \langle   u_{n-1} , v_1^{(2) (n-1)}, \ldots, v_{n-1}^{(2) (n-1)}  \rangle    \big) \triangleleft_{ \ll - , \cdots ,- \gg } (v_1^{(1)} \otimes \cdots \otimes v_{n-1}^{(1)}) \\
        & = \ll {\scriptstyle  \langle u_1, v_1^{(2)(1)}, \ldots, v_{n-1}^{(2) (1)} \rangle , v_{n-1}^{(1) (1)}, \ldots, v_{1}^{(1) (1)} } \gg \otimes \cdots \otimes \ll {\scriptstyle \langle u_{n-1}, v_1^{(2)(n-1)}, \ldots, v_{n-1}^{(2) (n-1)} \rangle , v_{n-1}^{(1) (n-1)}, \ldots, v_{1}^{(1) (n-1)} } \gg \\
         &= \big(   \varepsilon (v_1^{(1)}) \cdots \varepsilon ( v_{n-1}^{(1)}) u_1  \big) \otimes \cdots \otimes \big( \varepsilon (v_1^{(n-1)}) \cdots \varepsilon (v_1^{(n-1)}) u_{n-1}   \big) \\
         &= \big( \varepsilon (v_1^{(1)}) \cdots \varepsilon (v_1^{(n-1)})    \big) \cdots \big(   \varepsilon (v_{n-1}^{(1)}) \cdots \varepsilon (v_{n-1}^{(n-1)})   \big) (u_1 \otimes \cdots u_{n-1}) \\
        & = \varepsilon_{C^{\otimes (n-1)}} (v_1 \otimes \cdots \otimes v_{n-1}) (u_1 \otimes \cdots u_{n-1}).
     \end{align*}
     Similarly, we have
     \begin{align*}
        & \big( { (u_1 \otimes \cdots u_{n-1}) }\triangleleft_{\ll -, \ldots, - \gg} (v_1^{(2)} \otimes \cdots \otimes v_{n-1}^{(2)}) \big) \triangleleft_{ \langle - , \cdots ,- \rangle } (v_1^{(1)} \otimes \cdots \otimes v_{n-1}^{(1)}) \\
        & \qquad \quad = \varepsilon_{C^{\otimes (n-1)}} (v_1 \otimes \cdots \otimes v_{n-1}) ( { u_1 \otimes \cdots u_{n-1} }).
     \end{align*}
     This proves that $( C^{\otimes (n-1)} , \triangleleft_{\langle -, \ldots, - \rangle})$ is a linear rack.

     For the second part, we observe that $f : C \rightarrow C'$ is a coalgebra homomorphism implies that $f^{\otimes (n-1)} : C^{\otimes (n-1)} \rightarrow C'^{\otimes (n-1)}$ is also a coalgebra homomorphism. Also, we have
     \begin{align*}
        &f^{\otimes (n-1)} \big(    (u_1 \otimes \cdots \otimes u_{n-1}) \triangleleft_{\langle -, \ldots, - \rangle}  (v_1 \otimes \cdots \otimes v_{n-1})\big) \\
        %&= f (\langle u_1 , v_1^{(1)}, \ldots, v_{n-1}^{(1)}   \rangle ) \otimes \cdots \otimes f (\langle u_{n-1} , v_1^{(n-1)}, \ldots, v_{n-1}^{(n-1)} \rangle ) \\
        &= \langle  f(u_1), f(v_1^{(1)}), \ldots, f (v_{n-1}^{(1)})    \rangle' \otimes \cdots \otimes \langle  f (u_{n-1}), f (v_1^{(n-1)}), \ldots, f (v_{n-1}^{(n-1)}) \rangle' \\
        &= \langle  f(u_1), f(v_1)^{(1)}, \ldots, f (v_{n-1})^{(1)}    \rangle' \otimes \cdots \otimes \langle  f (u_{n-1}), f (v_1)^{(n-1)}, \ldots, f (v_{n-1})^{(n-1)} \rangle' \\
        &= (f^{\otimes (n-1)} (u_1 \otimes \cdots \otimes u_{n-1})) \triangleleft_{\langle -, \ldots, - \rangle'} (f^{\otimes (n-1)} (v_1 \otimes \cdots \otimes v_{n-1}))
     \end{align*}
     which shows that $f^{\otimes (n-1)} : C^{\otimes (n-1)} \rightarrow C'^{\otimes (n-1)}$ is a homomorphism of linear racks.
\end{proof}

\end{document}
